\documentclass[aspectratio=169,11pt]{beamer}
\makeatletter
\input{../../templates/dune-professional/beamerthemeDUNEProfessional.sty}
\makeatother
\DuneUseWhite
\title{What I added in PR \#8}
\subtitle{WIB, DAPHNE and PDU emulators for bridge tests}
\author{Manuel Arroyave}
\date{10 November 2026}
\hypersetup{colorlinks=true,urlcolor=DuneInfo,linkcolor=DuneInfo}

\begin{document}
\begin{frame}[plain]
  \titlepage
\end{frame}

\begin{frame}{Threads and asyncio handle waiting differently}
  \vfill
  {\normalsize
  \renewcommand{\arraystretch}{1.7}
  \begin{tabular}{@{}p{0.22\textwidth}p{0.35\textwidth}p{0.35\textwidth}@{}}
    & \textbf{Threads} & \textbf{Asyncio} \\
    \textbf{Switching} & Operating system & Event loop, using \texttt{await} \\
    \textbf{Slow device} & Other threads continue & Other tasks continue \\
    \textbf{Library calls} & Can use blocking calls & Needs async calls \\
    \textbf{Watch for} & Shared data may need locks & A blocking call can stall the loop
  \end{tabular}\par}
  \vfill
  {\small A blocking WIB read can hold up the loop until its timeout.\par
    Use an async read or a worker thread.\par}
\end{frame}

\begin{frame}{1. I added three simulated devices}
  \framesubtitle{\href{https://github.com/DUNE/dune-dscs/pull/8}{DUNE/dune-dscs, PR \#8}}
  \vfill
  \begin{columns}[T,onlytextwidth]
    \begin{column}{0.30\textwidth}
      {\LARGE\bfseries WIB\par}\medskip
      ZeroMQ + protobuf
    \end{column}
    \begin{column}{0.30\textwidth}
      {\LARGE\bfseries DAPHNE\par}\medskip
      ZeroMQ + protobuf v2
    \end{column}
    \begin{column}{0.30\textwidth}
      {\LARGE\bfseries PDU\par}\medskip
      SNMP v2c / Raritan
    \end{column}
  \end{columns}
  \bigskip\bigskip
  {\Large A bridge can talk to them without the hardware.\par}
  \vfill
\end{frame}

\begin{frame}{2. I used the existing message definitions}
  \vfill
  {\Large
    WIB and DAPHNE: bundled protobuf definitions.\par\bigskip
    PDU: Raritan PDU2-MIB outlet commands.\par\bigskip
    I recorded the source revisions and hashes.\par}
  \vfill
\end{frame}

\begin{frame}{3. I made commands change simulated state}
  \vfill
  {\large
  \renewcommand{\arraystretch}{1.7}
  \begin{tabular}{@{}p{0.22\textwidth}p{0.72\textwidth}@{}}
    \textbf{WIB} & FEMB power, registers and timing state \\
    \textbf{DAPHNE} & Configuration and trigger thresholds \\
    \textbf{PDU} & Outlet on/off state
  \end{tabular}\par}
  \bigskip\bigskip
  {\Large A write can be checked with a read.\par}
  \vfill
\end{frame}

\begin{frame}{4. I made each device easy to run}
  \vfill
  {\normalsize\texttt{uv run python -m dune\_dscs.emulator wib}\par}
  \bigskip\bigskip
  {\Large
    Choose \texttt{wib}, \texttt{daphne} or \texttt{pdu}.\par\bigskip
    Each runs in its own process, using asyncio.\par}
  \vfill
  {\small Host and port are configurable; defaults use loopback.\par}
\end{frame}

\begin{frame}{5. I checked the emulators}
  \framesubtitle{Local tests and Ignition on dune-fd-test01}
  \vfill
  {\Large
    Ran 963 devices locally.\par\bigskip
    Read 627 devices in Ignition.\par\bigskip
    Checked writes, device loss and recovery.\par}
  \vfill
  {\small The Ignition test used a separate adapter, outside PR \#8.\par
    Hardware and crate IPMI still need testing.\par}
\end{frame}

\begin{frame}{6. Geoff can use them to test his bridge}
  \vfill
  {\Large
    Point the bridge's device clients at the emulators.\par\bigskip
    Publish the results over OPC UA to Ignition.\par\bigskip
    Start with one device, then increase the fleet.\par}
  \vfill
  {\small PR \#8 supplies the emulators. The new bridge is still to be written.\par}
\end{frame}
\end{document}
